\section{Diagnostics: density sensitivity and diagram trajectories}
\label{app:diagnostics}

% Advisor round-2 items: M-sensitivity table + persistence-diagram
% trajectory visualization. All numbers regenerable:
% src: experiments/density_bound.py-style sweep (M in {512,1024,2048,4096},
% seed 0, floor = 6*r_med(M)); figure from the committed run dumps
% (sphere_{C0,C1_r0.1,C2_r0.1}_s0 at steps {0,400,1200,2500}).

\paragraph{Sampling density vs.\ measurability.}
Table~\ref{tab:msens} counts the significant bars of the clean target
surfaces across sampling densities under the self-calibrating floor
$6\,r_{\mathrm{med}}(M)$ of \S\ref{sec:discussion}. It confirms the three
density claims made in the paper: the sphere's void clears the floor at
every tested density; the torus needs $M{=}2048$ before its \emph{second}
loop is measurable, and its own void becomes measurable only at $M{=}4096$
(hence the H1 restriction at $M{=}2048$); and the double torus's tube
loops cross the floor between $M{=}2048$ and $M{=}4096$ --- consistent
with the $M\gtrsim3{\times}10^{3}$ estimate of \S\ref{sec:discussion} ---
so all four loops are measurable at $M{=}4096$.

\begin{table}[h]
\centering\small
\caption{Significant bars (H1 / H2) on the clean target surfaces vs.\
sampling density $M$; floor $=6\,r_{\mathrm{med}}(M)$, seed 0.}
\label{tab:msens}
\begin{tabular}{lcccc}
\toprule
 & $M{=}512$ & $M{=}1024$ & $M{=}2048$ & $M{=}4096$ \\
\midrule
% src: density sweep 2026-07-09 (floors .1261/.0895/.0644/.0458 sphere;
% .1228/.0843/.0623/.0437 torus; .1107/.0739/.0539/.0382 double_torus)
sphere       & 0 / 1 & 0 / 1 & 0 / 1 & 0 / 1 \\
torus        & 1 / 0 & 1 / 0 & 2 / 0 & 2 / 1 \\
double torus & 0 / 0 & 2 / 0 & 2 / 0 & 4 / 0 \\
\bottomrule
\end{tabular}
\end{table}

\paragraph{Diagram trajectories: the pinning mechanism, visualized.}
Fig.~\ref{fig:pdtraj} shows the live H2 diagrams of one sphere seed under
C0/C1/C2 at four training steps. All three arms start with the void
essentially on target (initialization samples the SfM point cloud), and
photometric training alone drags it away (the baseline settles
${\approx}.06$ below). The loss arm dips identically while the ramp is
inactive, then pulls the void back to within ${\approx}.016$ of target ---
matching its measured tail. The control arm is qualitatively different:
the void's death is dragged progressively down (${\approx}.23 \to .15 \to
.13$) while a carpet of phantom mid-persistence bars accumulates near the
diagonal --- the repulsion-equilibrium signature discussed in
\S\ref{sec:discussion}.

\begin{figure}[tb]
\centering
% src: figures/pd_trajectories.png (generated from committed dumps)
\includegraphics[width=\linewidth]{pd_trajectories.png}
\caption{Live H2 persistence diagrams over training (one sphere seed; 20k
samples, target-normalized; star $=$ target void bar, circle $=$ most
persistent live bar). \emph{Top}, baseline: the void drifts off target.
\emph{Middle}, loss: the void returns to the target once the ramp
activates. \emph{Bottom}, repulsion control: the void is pinned and
dragged away while near-diagonal phantom bars accumulate.}
\label{fig:pdtraj}
\end{figure}
